% SPDX-License-Identifier: MIT-0
\documentclass[11pt,a4paper]{article}
\usepackage[T1]{fontenc}
\usepackage[utf8]{inputenc}
\usepackage{lmodern}
\usepackage[margin=25mm,headheight=14pt]{geometry}
\usepackage{amsmath,amssymb,amsthm,mathtools}
\usepackage{booktabs,array,longtable}
\usepackage{microtype}
\usepackage[dvipsnames]{xcolor}
\usepackage{fancyhdr}
\usepackage{enumitem}
\usepackage{listings}
\usepackage{titlesec}
\usepackage[colorlinks=true,linkcolor=MidnightBlue,citecolor=MidnightBlue,
            urlcolor=MidnightBlue,breaklinks=true]{hyperref}
\definecolor{Ink}{HTML}{20354C}
\definecolor{Accent}{HTML}{536147}
\titleformat{\section}{\large\bfseries\color{Ink}}{\thesection}{0.7em}{}
\titleformat{\subsection}{\normalsize\bfseries\color{Accent}}{\thesubsection}{0.7em}{}
\pagestyle{fancy}
\fancyhf{}
\fancyhead[L]{\small Stirling subset polynomials and Hankel obstructions}
\fancyhead[R]{\small Research draft}
\fancyfoot[C]{\thepage}
\renewcommand{\headrulewidth}{0.25pt}
\newtheorem{theorem}{Theorem}[section]
\newtheorem{lemma}[theorem]{Lemma}
\newtheorem{proposition}[theorem]{Proposition}
\newtheorem{corollary}[theorem]{Corollary}
\theoremstyle{definition}
\newtheorem{definition}[theorem]{Definition}
\newtheorem{remark}[theorem]{Remark}
\numberwithin{equation}{section}
\newcommand{\N}{\mathbb N}
\newcommand{\R}{\mathbb R}
\newcommand{\Z}{\mathbb Z}
\newcommand{\coeff}[2]{[#1^{#2}]}
\newcommand{\St}[2]{S^{(r)}_{#1,#2}}
\newcommand{\sr}[1]{s_{r,#1}(x)}
\newcommand{\HT}{H_{r,a}(x)}
\newcommand{\dd}{\mathop{}\!\mathrm{d}}
\DeclareMathOperator{\ord}{ord}
\DeclareMathOperator{\sgn}{sgn}
\emergencystretch=2em
\setlength{\parskip}{0.25em}
\setlist{nosep}
\lstset{basicstyle=\small\ttfamily,breaklines=true,columns=fullflexible,
        frame=single,rulecolor=\color{black!20},xleftmargin=0.3em,xrightmargin=0.3em}
\hypersetup{pdftitle={A uniform Hankel obstruction for higher-order Stirling subset polynomials},
 pdfauthor={AI-assisted research draft prepared for Vladimir Reshetnikov},
 pdfsubject={A proof of the higher-order negative assertion in Deb--Sokal Conjecture 1.4(d)}}

\begin{document}
\begin{center}
{\LARGE\bfseries\color{Ink} A uniform Hankel obstruction for\\[0.18em]
higher-order Stirling subset polynomials\par}
\vspace{0.8em}
{\large A proof of the higher-order assertion in Deb--Sokal Conjecture 1.4(d)}\\[0.9em]
{Prepared for Vladimir Reshetnikov}\\[0.3em]
{20 September 2026}
\end{center}

\begin{abstract}
Let $S^{(r)}_{n,k}$ count partitions of $n+(r-1)k$ labels into $k$ blocks,
each of size at least $r$, and put $s_{r,n}(x)=\sum_k S^{(r)}_{n,k}x^k$.
For every integer $r\ge3$ we prove that this polynomial sequence is not
coefficientwise Hankel-totally positive. More strongly, every adjacent
$3\times3$ Hankel determinant with positive shift has a negative coefficient
of $x^5$, and all its lower coefficients vanish. For shift one, the coefficient is
\[
-\binom{2r-1}{r-1}^{\!2}
 \frac{(r-2)(r-1)(r^2+11r+6)}{(r+1)(r+2)^2(r+3)}.
\]
The proof reduces a rank-one perturbation determinant to the second
differences of the two-block counting sequence. Those differences are a
geometric sequence minus a partial binomial sum; an elementary ratio
argument determines the sign uniformly. This proves the previously
conjectured negative range of Deb--Sokal Conjecture~1.4(d), using the
shift-one witness already suggested by its authors. We also obtain exact
formulas for arbitrary shifts, quantitative negative bounds, small-parameter
obstructions to both Stieltjes and Hamburger moment representations, and
asymptotic formulas. The accompanying code supplies independent
exact-arithmetic checks, not a replacement for the proofs.
\end{abstract}

\medskip
\noindent\textbf{Status.}
This is an AI-assisted research draft. The negative-range proofs are
self-contained; the known positive cases $r=1,2$ are imported when stating
the complete dichotomy. No independent referee or proof-assistant check is
claimed. The source audit did not locate an intervening proof, but it does
not establish priority.

\setcounter{tocdepth}{1}
\makeatletter
\renewcommand*\l@section{\@dottedtocline{1}{0em}{1.5em}}
\makeatother
\tableofcontents
\clearpage

\section{The problem and its relation to the supplied manifest}

Deb and Sokal's Conjecture~1.4(d) asks for a sharp cutoff in coefficientwise
Hankel positivity of higher-order Stirling subset polynomials: positivity
at orders one and two, and failure at every higher order
\cite[Conjecture~1.4(d), p.~10]{DebSokal}. Their Appendix~C suggests the
shift-one $3\times3$ determinant as a witness, without an all-order sign
formula \cite[p.~50]{DebSokal}. The present argument proves that prediction
and extends it to every positive shift.

The supplied manifest contains an entry entitled \emph{A sharp order-five
threshold for higher-order Stirling subset log-concavity}
\cite{Manifest}. That entry concerns part~(c) of the same conjecture:
log-concavity \emph{along individual rows} of the counting triangle. Here
we study part~(d): determinants built \emph{across different row-generating
polynomials}. The underlying numbers are the same, but the assertions are
not duplicates. No theorem claimed by that manifest entry is used below.
The manifest is treated as a catalogue of reported results, not as an
independent verification of their proofs.

The useful feature of the selected problem is that a uniform negative
answer can be certified by one coefficient of one fixed-size minor.
There is no need to search larger and larger matrices or to control every
coefficient of the row polynomials. In fact, the coefficient of $x^2$ in
each row already contains all the information needed to obtain the
obstruction.

\section{Definitions and principal results}
\label{sec:definitions}

\subsection{The counting triangle}

For integers $r\ge1$ and $0\le k\le n$, let $S^{(r)}_{n,k}$ be the number
of set partitions of an $(n+(r-1)k)$-element labelled set into exactly $k$
nonempty blocks, with every block having size at least $r$. The empty
partition gives $S^{(r)}_{0,0}=1$. Entries outside the triangle are zero.
Define
\begin{equation}
 s_{r,n}(x)=\sum_{k=0}^{n}S^{(r)}_{n,k}x^k.
 \label{eq:rows}
\end{equation}
In particular, $s_{r,0}(x)=1$, while for every $n\ge1$,
\begin{equation}
 s_{r,n}(x)=x+b_nx^2+O(x^3),\qquad b_n=S^{(r)}_{n,2}.
 \label{eq:smallx}
\end{equation}
Here and throughout, $O(x^j)$ can be read formally as divisibility by
$x^j$ in a power-series ring. No analytic convergence is involved in
coefficient extraction. The coefficient of $x$ is one because there is
exactly one partition into a single block of the required size.

For completeness, the triangular recurrence is
\begin{equation}
 S^{(r)}_{n,k}=
 \binom{n+(r-1)k-1}{r-1}S^{(r)}_{n-1,k-1}
       +kS^{(r)}_{n-1,k}.
 \label{eq:triangle-recurrence}
\end{equation}
To see this, distinguish the largest label. If its block has size exactly
$r$, choose the other $r-1$ labels and remove that block. If its block has
size greater than $r$, remove only the distinguished label, then record
which of the $k$ blocks receives it on reinsertion. Both cases leave
precisely the label counts encoded on the right-hand side.

\subsection{Hankel positivity and the witness}

\begin{definition}
A polynomial sequence $(p_n(x))_{n\ge0}$ is \emph{coefficientwise
Hankel-totally positive} when every finite minor of
$(p_{i+j}(x))_{i,j\ge0}$ has only nonnegative coefficients.
\end{definition}

We follow the source's convention that ``totally positive'' allows zero
minors; this is also called total nonnegativity. In particular, finding
one strictly negative coefficient disproves the property.

For a positive integer $a$, put
\begin{equation}
 H_{r,a}(x)=\det\bigl(s_{r,a+i+j}(x)\bigr)_{0\le i,j\le2}.
 \label{eq:H-def}
\end{equation}
This is a minor of the full Hankel matrix: take row indices $a,a+1,a+2$
and column indices $0,1,2$. Let
\begin{equation}
 K_{r,a}=[x^5]H_{r,a}(x),\qquad
 F_m(N)=\sum_{j=0}^{m}\binom Nj,
 \label{eq:K-F}
\end{equation}
with $F_{-1}(N)=F_{-2}(N)=0$ and $\binom N{-1}=0$.

\begin{theorem}[Uniform shifted obstruction]
\label{thm:shifted}
Let $r\ge3$ and $a\ge1$ be integers. Set $N=a+2r-2$ and $m=r-3$. Then
\begin{align}
 H_{r,a}(x)&=K_{r,a}x^5+O(x^6),\label{eq:leading}\\
 K_{r,a}&=-2^{N-1}\left(\binom Nm-\binom N{m-1}\right)
        +F_m(N)F_m(N+2)-F_m(N+1)^2.\label{eq:general-K}
\end{align}
Moreover,
\begin{equation}
 K_{r,a}\le
 -2^{N-1}\left(\binom Nm-\binom N{m-1}\right)<0.
 \label{eq:negative-bound}
\end{equation}
Consequently every determinant in \eqref{eq:H-def} has order of vanishing
exactly five at the origin and a negative leading coefficient.
\end{theorem}

\begin{theorem}[Factored first-shift certificate]
\label{thm:factored}
For every integer $r\ge1$,
\begin{equation}
 \boxed{\displaystyle
 K_{r,1}=-\binom{2r-1}{r-1}^{\!2}
 \frac{(r-2)(r-1)(r^2+11r+6)}{(r+1)(r+2)^2(r+3)}.}
 \label{eq:factored}
\end{equation}
Thus $K_{r,1}=0$ for $r=1,2$ and $K_{r,1}<0$ for every $r\ge3$.
\end{theorem}

Theorem~\ref{thm:factored} alone proves the new negative-range assertion.
Theorem~\ref{thm:shifted} explains its mechanism and gives additional
consequences. These are two compatible routes to the same sign, not
independent evidence from a finite list of examples.

\begin{corollary}[The cutoff]
\label{cor:cutoff}
For positive integer $r$, the sequence $(s_{r,n}(x))_{n\ge0}$ is
coefficientwise Hankel-totally positive if and only if $r\in\{1,2\}$.
\end{corollary}
\begin{proof}
The negative range follows from either theorem. The positive cases are
known results: the ordinary Stirling case and the Ward-polynomial case,
respectively. The latter is supported by the continued-fraction result of
Elvey Price and Sokal \cite{ElveyPriceSokal}; the deduction of the required
positivity, and the status of both positive cases, are recorded explicitly
in \cite[p.~10]{DebSokal}. We do not count these positive cases as new
results of this article.
\end{proof}

\section{A rank-one reduction for the leading coefficient}
\label{sec:rankone}

The determinant is initially degenerate to unusually high order: after
factoring one $x$ from each row, its constant matrix is the all-ones
matrix, which has rank one. Adjacent differences isolate the first
nontrivial part.

\begin{lemma}[Low-degree Hankel reduction]
\label{lem:rankone}
Suppose $p_j(x)=x+b_jx^2+O(x^3)$ for the five indices
$j=a,\ldots,a+4$. With $\Delta b_j=b_{j+1}-b_j$,
\begin{equation}
 \det(p_{a+i+j}(x))_{0\le i,j\le2}
 =x^5\left((\Delta^2b_a)(\Delta^2b_{a+2})
                -(\Delta^2b_{a+1})^2\right)+O(x^6).
 \label{eq:rankone-identity}
\end{equation}
In particular all coefficients below degree five vanish, and coefficients
of $p_j$ of degree at least three cannot affect the displayed term.
\end{lemma}
\begin{proof}
Factor $x$ from each row. The remaining matrix is
$J+xB+O(x^2)$, where $J$ is the all-ones matrix and
$B_{ij}=b_{a+i+j}$. The matrix
\[
 L=\begin{pmatrix}1&0&0\\-1&1&0\\0&-1&1\end{pmatrix}
\]
has determinant one. Multiplication on the left by $L$ and on the right
by $L^{\mathsf T}$ preserves the determinant. Since
$LJL^{\mathsf T}=\operatorname{diag}(1,0,0)$, the resulting matrix has the
block form
\begin{equation}
 \begin{pmatrix}
  1+O(x)&O(x)&O(x)\\
  O(x)&x\Delta^2b_a+O(x^2)&x\Delta^2b_{a+1}+O(x^2)\\
  O(x)&x\Delta^2b_{a+1}+O(x^2)&x\Delta^2b_{a+2}+O(x^2)
 \end{pmatrix}.
 \label{eq:transformed}
\end{equation}
The terms of degree two in its determinant are exactly the determinant
of the displayed lower-right block of linear terms. All permutations
using an off-diagonal entry in the first row also use an off-diagonal
entry in the first column, and hence have degree at least three. Restoring
the factor $x^3$ proves \eqref{eq:rankone-identity}.
\end{proof}

The sign question is therefore equivalent to a strict log-concavity
inequality for the sequence $d_n=\Delta^2b_n$. This is log-concavity in the
\emph{row index after taking two differences}, not the log-concavity of
$k\mapsto S^{(r)}_{n,k}$ considered in the manifest.

\section{Two-block partitions and the all-shifts proof}
\label{sec:allshifts}

\subsection{The two-block column}

\begin{lemma}
\label{lem:two-blocks}
For $n\ge1$ and $N=n+2r-2$,
\begin{equation}
 b_n=S^{(r)}_{n,2}
 =\frac12\sum_{j=r}^{N-r}\binom Nj
 =2^{N-1}-\sum_{j=0}^{r-1}\binom Nj.
 \label{eq:b-formula}
\end{equation}
The first sum is empty when $n=1$, giving $b_1=0$.
For $r\ge3$,
\begin{equation}
 d_n:=\Delta^2b_n=2^{N-1}-F_{r-3}(N).
 \label{eq:d-formula}
\end{equation}
For $r=1,2$, instead $d_n=2^{N-1}$.
\end{lemma}
\begin{proof}
An unordered partition into two admissible blocks is specified twice by
choosing one block, whose size must lie between $r$ and $N-r$.
The omitted lower and upper binomial tails are equal by symmetry. Since
$N\ge2r-1$, the tails are disjoint, including at the boundary $n=1$.
Subtracting them from $2^N$ gives \eqref{eq:b-formula}.

When $n$ increases by one, so does $N$. Pascal's identity gives
\[
 \Delta\binom Nj=\binom N{j-1},\qquad
 \Delta^2\binom Nj=\binom N{j-2},
\]
where terms with negative lower index are zero. Also
$\Delta^2 2^{N-1}=2^{N-1}$. Applying these identities term by term gives
\eqref{eq:d-formula}; the finite polynomial correction disappears
entirely when $r=1,2$.
\end{proof}

For $r\ge3$, the relevant $N$ satisfies $N\ge2(r-3)+5$.
Binomial symmetry then gives $F_{r-3}(N)<2^{N-1}$, since its final
index lies strictly below the midpoint and omits central terms. Thus
$d_n>0$; the sequence whose strict log-concavity we use is positive.

\subsection{A partial-binomial log-concavity lemma}

\begin{lemma}
\label{lem:F-logconcave}
For fixed integer $m\ge0$, the positive sequence $F_m(N)$, $N\ge m$,
is log-concave:
\begin{equation}
 F_m(N)F_m(N+2)-F_m(N+1)^2\le0.
 \label{eq:F-logconcave}
\end{equation}
Equality holds for $m=0$, and the inequality is strict for $m\ge1$.
\end{lemma}
\begin{proof}
The case $m=0$ is $F_0(N)=1$. Suppose $m\ge1$. Pascal's identity gives
\begin{equation}
 \frac{F_m(N+1)}{F_m(N)}
 =1+\frac{F_{m-1}(N)}{F_m(N)}.
 \label{eq:F-ratio}
\end{equation}
For each $0\le j<m$,
\begin{equation}
 \frac{\binom Nj}{\binom Nm}
 =\frac{m!}{j!}\,
   \frac{1}{\prod_{h=j}^{m-1}(N-h)}
 \label{eq:binom-ratio}
\end{equation}
is strictly decreasing with integer $N\ge m$. Hence so is the positive
sum
\[
 R_m(N)=\frac{F_{m-1}(N)}{\binom Nm}.
\]
The ratio appearing in \eqref{eq:F-ratio} is
$F_{m-1}(N)/F_m(N)=R_m(N)/(1+R_m(N))$, a strictly increasing function of
$R_m(N)$. It therefore decreases with $N$. Successive ratios of $F_m$
strictly decrease, which is exactly \eqref{eq:F-logconcave} after
multiplication by the positive denominators.
\end{proof}

This proof deliberately uses only finite sums and positive ratios. It
requires no general preservation theorem for log-concavity.

\subsection{Proof of Theorem \ref{thm:shifted}}

Set $B=2^{N-1}$ and abbreviate $F_j=F_m(N+j)$ for $j=0,1,2$ in this
subsection only. By Lemmas~\ref{lem:rankone} and~\ref{lem:two-blocks},
\begin{align}
 K_{r,a}
   &=(B-F_0)(4B-F_2)-(2B-F_1)^2\notag\\
   &=-B(F_2-4F_1+4F_0)+F_0F_2-F_1^2.
 \label{eq:K-expand}
\end{align}
The $B^2$ terms cancel. Applying Pascal twice yields
\begin{align*}
 F_m(N+1)&=F_m(N)+F_{m-1}(N),\\
 F_m(N+2)&=F_m(N)+2F_{m-1}(N)+F_{m-2}(N).
\end{align*}
Consequently
\begin{align}
 F_m(N+2)-4F_m(N+1)+4F_m(N)
  &=F_m(N)-2F_{m-1}(N)+F_{m-2}(N)\notag\\
  &=\binom Nm-\binom N{m-1}.
 \label{eq:tail-cancellation}
\end{align}
This proves \eqref{eq:general-K}.

Here $N=a+2r-2\ge2r-1=2m+5$. If $m=0$, the binomial difference in
\eqref{eq:tail-cancellation} is one. If $m\ge1$, it is positive because
\[
 \frac{\binom Nm}{\binom N{m-1}}=\frac{N-m+1}{m}>1.
\]
Lemma~\ref{lem:F-logconcave} makes the remaining term of
\eqref{eq:K-expand} nonpositive. This proves the strict negativity and
the quantitative bound. The rank-one lemma proves the vanishing of the
lower coefficients, so the order of vanishing is exactly five.
\hfill$\square$

\begin{remark}[The two lowest orders]
For $r=1,2$, Lemma~\ref{lem:two-blocks} makes $d_n$ geometric. Its adjacent
$2\times2$ determinants vanish, and hence $K_{r,a}=0$ for every $a\ge1$.
This explains why the degree-five obstruction starts at $r=3$.
Vanishing at one degree is not itself a proof of total positivity at
$r=1,2$; those cases remain the imported positive results of
Corollary~\ref{cor:cutoff}.
\end{remark}

\section{A short factored certificate at shift one}
\label{sec:factor}

We now give a separate, very small algebraic proof of
Theorem~\ref{thm:factored}. It uses only the first five entries of the
two-block column and is useful as a standalone certificate.

Let
\[
 A=\binom{2r-1}{r-1},\qquad
 u=\frac{2r}{r+1},\qquad
 v=\frac{2r(2r+1)}{(r+1)(r+2)},\qquad
 w=\frac{2r(2r+1)(2r+2)}{(r+1)(r+2)(r+3)}.
\]
The recurrence \eqref{eq:triangle-recurrence} at $k=2$ is
\begin{equation}
 b_n=2b_{n-1}+\binom{n+2r-3}{r-1}\qquad(n\ge2).
 \label{eq:b-recur}
\end{equation}
The successive binomial terms at $n=2,3,4,5$ are $A,Au,Av,Aw$.
Starting from $b_1=0$ gives
\begin{equation}
 (b_1,b_2,b_3,b_4,b_5)
 =A\bigl(0,1,2+u,4+2u+v,8+4u+2v+w\bigr).
 \label{eq:first-b}
\end{equation}
Thus
\begin{equation}
 (d_1,d_2,d_3)=A\bigl(u,1+v,2+u+w\bigr).
 \label{eq:first-d}
\end{equation}
The rank-one identity immediately gives
\begin{equation}
 \frac{K_{r,1}}{A^2}=u(2+u+w)-(1+v)^2.
 \label{eq:short-certificate}
\end{equation}
For a completely explicit factorization, observe that
\[
 1+v=\frac{5r^2+5r+2}{(r+1)(r+2)},\qquad
 2+u+w=\frac{2(2r+1)(3r^2+7r+6)}{(r+1)(r+2)(r+3)}.
\]
On the common denominator $(r+1)^2(r+2)^2(r+3)$, the numerator in
\eqref{eq:short-certificate} becomes
\begin{align*}
 &4r(2r+1)(3r^2+7r+6)(r+2)-(5r^2+5r+2)^2(r+3)\\
 &\hspace{2em}=-(r-2)(r-1)(r+1)(r^2+11r+6).
\end{align*}
Canceling the factor $r+1$ proves \eqref{eq:factored}. All denominators
are positive for $r\ge1$, and every factor in the numerator after the
minus sign is positive for $r\ge3$. No interpolation in $r$ is being
used: this is an identity of rational functions.

\begin{table}[ht]
\centering
\caption{The first-shift degree-five coefficient. Every entry is an exact integer.}
\label{tab:K}
\begin{tabular}{@{}rrr@{}}
\toprule
$r$ & $\binom{2r-1}{r-1}$ & $K_{r,1}$\\
\midrule
1&1&0\\
2&3&0\\
3&10&$-16$\\
4&35&$-385$\\
5&126&$-6\,966$\\
6&462&$-114\,345$\\
7&1\,716&$-1\,799\,512$\\
8&6\,435&$-27\,756\,729$\\
9&24\,310&$-423\,939\,880$\\
10&92\,378&$-6\,445\,028\,304$\\
\bottomrule
\end{tabular}
\end{table}

\section{Worked examples and moment obstructions}
\label{sec:examples}

\subsection{The smallest higher order}

For $r=3$, recurrence \eqref{eq:triangle-recurrence} gives
\begin{align*}
 s_{3,1}(x)&=x,\\
 s_{3,2}(x)&=x+10x^2,\\
 s_{3,3}(x)&=x+35x^2+280x^3,\\
 s_{3,4}(x)&=x+91x^2+2100x^3+15400x^4,\\
 s_{3,5}(x)&=x+210x^2+10395x^3+200200x^4+1401400x^5.
\end{align*}
Their determinant is
\begin{equation}
 \begin{split}
 H_{3,1}(x)=x^5\bigl(&-16+19110x+938700x^2\\
                    &+15309000x^3+79380000x^4\bigr).
 \end{split}
 \label{eq:H31}
\end{equation}
The entire negative coefficient can also be checked without expanding
these polynomials: $(b_1,\ldots,b_5)=(0,10,35,91,210)$ has second
differences $(15,31,63)$, and $15\cdot63-31^2=-16$.

For comparison, the two lower-order determinants are
\begin{align}
 H_{1,1}(x)&=2x^6,\notag\\
 H_{2,1}(x)&=x^6(216+960x+1440x^2+720x^3).
 \label{eq:low-order-H}
\end{align}
These examples illustrate the threshold but are not used to extrapolate
it to other orders.

In \eqref{eq:H31}, the quartic after $x^5$ is strictly increasing for
$x>0$, starts negative, and tends to positive infinity. It has exactly
one positive root $\xi$, and exact rational evaluations give
\begin{equation}
 0.000805006270<\xi<0.000805006271.
 \label{eq:root-bracket}
\end{equation}
In particular,
\begin{equation}
 H_{3,1}(1/2000)=
 -\frac{4966725131}{25600000000000000000000000}<0.
 \label{eq:scalar-example}
\end{equation}
The bracket is only an illustration. The proof for all $r$ and all shifts
requires no numerical root finding.

\subsection{Failure of scalar moment representations}

A negative polynomial coefficient does not, by itself, force the
polynomial to be negative at a positive argument. Here it does force
negativity \emph{near zero}, because the negative coefficient is the
first nonzero one. This distinction is essential.

\begin{proposition}[An effective negative interval]
\label{prop:epsilon}
For $r\ge3$ and $a\ge1$, define $L_j=s_{r,j}(1)$ and
\begin{equation}
 \begin{split}
 M_{r,a}={}&L_aL_{a+2}L_{a+4}+L_aL_{a+3}^{\,2}
            +L_{a+1}^{\,2}L_{a+4}\\
           &+2L_{a+1}L_{a+2}L_{a+3}+L_{a+2}^{\,3},\\
 \varepsilon_{r,a}={}&\min\left(1,\frac{-K_{r,a}}{2M_{r,a}}\right).
 \end{split}
 \label{eq:epsilon}
\end{equation}
Then $\varepsilon_{r,a}$ is a positive rational number and
\begin{equation}
 H_{r,a}(x)\le \tfrac12K_{r,a}x^5<0
 \qquad(0<x\le\varepsilon_{r,a}).
 \label{eq:epsilon-result}
\end{equation}
\end{proposition}
\begin{proof}
The five products defining $M_{r,a}$ are the unsigned terms in the
symmetric $3\times3$ determinant expansion, with the paired terms counted
twice. Since each row polynomial has nonnegative coefficients, the sum
of the absolute values of all coefficients of $H_{r,a}$ is at most
$M_{r,a}$. For $0<x\le1$ the terms of degree at least six consequently
have absolute sum at most $M_{r,a}x^6$. Use
$K_{r,a}<0$ and the definition of $\varepsilon_{r,a}$.
\end{proof}

\begin{corollary}[Stieltjes and Hamburger exclusions]
\label{cor:moments}
For every $r\ge3$, the following assertions hold.
For $0<x\le\varepsilon_{r,1}$, the numerical sequence
$(s_{r,n}(x))_{n\ge0}$ is not a Stieltjes moment sequence. For
$0<x\le\varepsilon_{r,2}$, it is not even a Hamburger moment sequence.
\end{corollary}
\begin{proof}
A Stieltjes representation would give a positive measure $\mu$ supported
on $[0,\infty)$ with
$s_{r,n}(x)=\int t^n\dd\mu(t)$. For every real vector
$v=(v_0,v_1,v_2)$,
\[
 \sum_{i,j=0}^2v_iv_j s_{r,a+i+j}(x)
  =\int t^a(v_0+v_1t+v_2t^2)^2\dd\mu(t)\ge0.
\]
Thus the matrix defining $H_{r,a}$ would be positive semidefinite and its
determinant nonnegative. Shift $a=1$ contradicts
Proposition~\ref{prop:epsilon}.

For a Hamburger representation the positive measure may be supported
anywhere on $\R$. Use $a=2$, for which the integrand is again
nonnegative. Equivalently, the matrix defining $H_{r,2}$ is the Gram
matrix of $t,t^2,t^3$ and a principal submatrix of the unshifted Hankel
matrix. Its negative determinant is again impossible.
\end{proof}

This argument uses only the necessary Gram-matrix condition for a
moment representation; no converse characterization theorem is needed.
The exclusions concern sufficiently small positive $x$. They do not
classify the numerical sequence at every positive specialization.

For a concrete Hamburger exclusion at $r=3$, the next shifted
determinant is
\begin{equation}
 \begin{split}
 H_{3,2}(x)=x^5\bigl(&-32+124614x+12875310x^2+556088400x^3\\
 &+12832394400x^4+165904200000x^5\\
 &+1136880360000x^6+3202824240000x^7\bigr).
 \end{split}
 \label{eq:H32}
\end{equation}
The parenthesized polynomial has positive nonconstant coefficients, and
its value at $x=10^{-4}$ is less than $-19$, as can be checked by rational
arithmetic. It is therefore negative throughout $0<x\le10^{-4}$.
Hence the third-order sequence has no Hamburger moment representation
throughout that explicit interval.

\section{Size, asymptotics, and useful specializations}
\label{sec:asymptotics}

\subsection{Low-order formulas at every shift}

The general formula becomes particularly simple for the first two
higher orders. When $r=3$, $m=0$ and the partial-binomial determinant
vanishes. When $r=4$, $m=1$ and $F_1(N)=N+1$, whose adjacent determinant
is $-1$. Therefore
\begin{align}
 K_{3,a}&=-2^{a+3},\label{eq:r3-all}\\
 K_{4,a}&=-2^{a+5}(a+5)-1.\label{eq:r4-all}
\end{align}
These show that the obstruction is not confined to an initial corner of
the Hankel matrix. Moving the test arbitrarily far down the sequence
preserves the sign.

\subsection{Large shift at fixed order}

\begin{proposition}
\label{prop:shift-asymptotic}
For each fixed $r\ge3$, as $a\to\infty$,
\begin{equation}
 K_{r,a}\sim
 -\frac{2^{a+2r-3}}{(r-3)!}\,a^{r-3}.
 \label{eq:shift-asymptotic}
\end{equation}
\end{proposition}
\begin{proof}
Set $m=r-3$ and $N=a+2r-2$. For fixed $m\ge1$,
\[
 \binom Nm-\binom N{m-1}
   =\frac{N^m}{m!}+O(N^{m-1}).
\]
Each $F_m(N+j)$ is a polynomial in $N$ of degree $m$, so its determinant
in \eqref{eq:general-K} is at most polynomial in $N$. It is negligible
compared with the exponential term $2^{N-1}N^m/m!$. Since $N/a\to1$,
\eqref{eq:shift-asymptotic} follows. For $m=0$, use the exact formula
\eqref{eq:r3-all}.
\end{proof}

\subsection{Large order at shift one}

The rational factor in \eqref{eq:factored} satisfies the useful identity
\begin{equation}
 \begin{split}
 Q(r)&:=\frac{(r-2)(r-1)(r^2+11r+6)}{(r+1)(r+2)^2(r+3)}\\
     &=1-\frac{24r(2r+1)}{(r+1)(r+2)^2(r+3)}
       =1+O(r^{-2}).
 \end{split}
 \label{eq:Qasym}
\end{equation}
In particular $0<-K_{r,1}<\binom{2r-1}{r-1}^2$ for every $r\ge3$.
Since $\binom{2r-1}{r-1}=\tfrac12\binom{2r}{r}$, the elementary
central-binomial asymptotic, obtained from Stirling's formula, yields
\begin{equation}
 K_{r,1}\sim-\frac{16^r}{4\pi r}\qquad(r\to\infty).
 \label{eq:order-asymptotic}
\end{equation}
Thus the coefficient is not merely negative by a vanishingly small
margin: its absolute size grows exponentially in the order. This does
not imply a uniform interval of negative \emph{specialized}
determinants, because the higher coefficients also grow.

\section{Exact verification and reproducibility}
\label{sec:computation}

The article's proofs establish the results for all parameter values.
The supplied programs audit indexing, cancellation, integrality,
factorization, and examples. They use no fitted recurrence and no
floating-point sign tests.

\subsection{Independent counting and determinant calculations}

The standard-library program \texttt{code/verify.py} constructs the
triangle using \eqref{eq:triangle-recurrence}. An independent routine
counts partitions by the block containing label one. Writing $T_r(M,k)$
for the number of admissible partitions of $M$ labels into $k$ blocks,
that routine uses
\begin{equation}
 T_r(M,k)=\sum_{j=r}^{M-r(k-1)}
       \binom{M-1}{j-1}T_r(M-j,k-1),
 \label{eq:independent-count}
\end{equation}
with $T_r(0,0)=1$ and impossible cases zero. This conditions on the
\emph{entire distinguished block}, rather than deleting one label as in
\eqref{eq:triangle-recurrence}, and checks the label shift
$M=n+(r-1)k$ directly.

The determinant routine expands all six signed permutation products
with exact integer polynomial arithmetic. A degree-five truncation is
also used, but it still retains enough coefficients to detect the
cancellation of all terms below $x^5$ and to compare the full six-term
expansion with the two-block formula. It is not a computation of the
same closed formula under a different name.

\begin{table}[ht]
\centering
\caption{Executed verification coverage. Counts are finite audit cases.}
\label{tab:checks}
\small
\begin{tabular}{@{}p{5.1cm}p{6.7cm}r@{}}
\toprule
Check & Range & Cases\\
\midrule
Independent partition counts & $1\le r\le6$, $0\le n\le8$, $0\le k\le n$ &270\\
Full polynomial determinants & $1\le r\le12$, $1\le a\le8$ &96\\
Degree-five determinant expansions & $3\le r\le30$, $1\le a\le40$ &1120\\
First-shift factorization & $1\le r\le100$ &100\\
Partial-binomial log-concavity & $0\le m\le30$, $\max(1,m)\le N\le100$ &2665\\
\bottomrule
\end{tabular}
\end{table}

All checks in Table~\ref{tab:checks} passed. The output records every full
determinant, a grid of shifted coefficients, the first one hundred
first-shift coefficients, the exact scalar examples, and a rational
bracket for the root in \eqref{eq:root-bracket}. The optional script
\texttt{code/symbolic\_identity.py} independently verifies the rational
factorization \eqref{eq:short-certificate} using SymPy~1.14.0. That is a
symbolic algebra check, not a proof-assistant certification.

\subsection{Commands and file layout}

The core verifier requires Python~3.9 or later and no third-party
packages. From the extracted directory, run
\begin{lstlisting}
python code/verify.py
\end{lstlisting}
The optional symbolic check additionally requires SymPy:
\begin{lstlisting}
python code/symbolic_identity.py
\end{lstlisting}
The document uses ordinary pdfLaTeX packages and an embedded bibliography:
\begin{lstlisting}
pdflatex -halt-on-error article.tex
pdflatex -halt-on-error article.tex
\end{lstlisting}
The archive also contains a \texttt{Makefile}, \texttt{README.md},
\texttt{STATUS.md}, a source audit, and a proof-audit note. The data files
are written deterministically; numerical timings and machine-dependent
paths are deliberately excluded from the mathematical certificates.

\section{Scope and remaining questions}
\label{sec:scope}

The established result is the all-$r\ge3$ failure claimed in
Conjecture~1.4(d), strengthened to a leading negative coefficient in
\emph{every} positive-shift adjacent $3\times3$ minor. Together with the
previously known positive cases this gives the full dichotomy in that
part of the conjecture.

Several nearby questions are not settled here. We do not prove total
positivity of the numerical triangle $S^{(r)}$ itself, total positivity
of any row reversal, row log-concavity, or Hankel positivity of the
higher-order \emph{cycle} polynomials. We also do not claim that all
$2\times2$ polynomial Hankel minors are nonnegative; hence the present
$3\times3$ witness is not asserted to have minimum possible size.

For scalar specializations, an especially concrete continuation is to
classify the set of $x>0$ for which $(s_{r,n}(x))_{n\ge0}$ is a
Stieltjes or Hamburger moment sequence. The small-$x$ interval is now
excluded, but positivity of a single determinant past its first root
would not establish the required positivity of all other minors.
Another question is whether an analogue of the rank-one reduction
organizes leading coefficients of larger shifted minors into tractable
finite-difference determinants. These are directions beyond, not
conditions on, the proof above.

\appendix
\section{Indexing and boundary audit}

The most error-prone feature of this problem is that changing the number
of blocks changes the number of labels in the definition of a row.
For example, $S^{(r)}_{n,1}$ counts partitions of $n+r-1$ labels, whereas
$S^{(r)}_{n,2}$ counts partitions of $n+2r-2$ labels. Treating both as
partitions of the same set would give the wrong two-block column.

The shift $a$ in \eqref{eq:H-def} is the index of the top-left entry,
not the size of the determinant. The first witness uses polynomials
$s_{r,1}$ through $s_{r,5}$. The positive-shift hypothesis is substantive:
$s_{r,0}=1$ is not divisible by $x$, so the rank-one argument in the form
used here does not apply to shift zero.

For the all-shifts sign formula, the auxiliary parameters satisfy
$m=r-3\ge0$ and $N=a+2r-2\ge2m+5$. Thus every binomial coefficient in a
denominator ratio is positive when it is used. The conventions
$F_{-1}=F_{-2}=0$ and $\binom N{-1}=0$ handle $r=3,4$ without silently
invoking undefined lower indices. The cases $r=1,2$ are treated directly
by the finite-difference identity, not by extending $m$ to negative
values in an unproved formula.

Finally, \eqref{eq:factored} has rational-looking factors but its value
is an integer: it equals a coefficient of an integer polynomial by the
proof. No denominator cancellation hypothesis or generic-parameter
specialization is needed. In particular, zeros at $r=1,2$ are harmless,
and there is no division by a Hankel determinant anywhere in the
argument.

\section{Proof dependency map}

The negative-range proof has a short dependency chain. The counting
definition gives \eqref{eq:smallx} and the two-block sum
\eqref{eq:b-formula}. The rank-one lemma converts the desired coefficient
into $d_ad_{a+2}-d_{a+1}^2$. Pascal's identity computes $d_a$ and the
geometric-term cancellation. A ratio comparison of finite binomial
sums proves the sign. Every step is an explicit algebraic or counting
identity.

The first-shift route replaces the binomial-sum sign lemma with a
rational factorization of three second differences. The moment
corollaries add only the elementary positivity of an integral of a
square and an explicit coefficient-norm estimate. The asymptotic
statements add standard estimates for fixed-degree polynomials and the
central binomial coefficient. Only the positive half of the final
cutoff imports a nontrivial published positivity theorem.

\begin{thebibliography}{9}

\bibitem{DebSokal}
Bishal Deb and Alan D. Sokal,
\emph{Higher-order Stirling cycle and subset triangles: Total positivity,
continued fractions and real-rootedness},
arXiv:2507.18959v1, 2025.
Conjecture~1.4(d) and discussion: p.~10;
proposed $3\times3$ witness: Appendix~C, p.~50.
\url{https://arxiv.org/abs/2507.18959v1}.

\bibitem{ElveyPriceSokal}
Andrew Elvey Price and Alan D. Sokal,
\emph{Phylogenetic trees, augmented perfect matchings, and a Thron-type
continued fraction (T-fraction) for the Ward polynomials},
Electronic Journal of Combinatorics \textbf{27}(4) (2020), P4.6.
\url{https://doi.org/10.37236/9571};
\url{https://arxiv.org/abs/2001.01468}.

\bibitem{Manifest}
Vladimir Reshetnikov,
\emph{A manifest of docs/reports}, supplied file \texttt{manifest.tex},
20 September 2026.
Entry \emph{A sharp order-five threshold for higher-order Stirling subset
log-concavity}, in the section on log-concavity, unimodality and
positivity. Used to identify a related, nonduplicate problem; not used
as proof evidence.

\end{thebibliography}
\end{document}
